\documentclass[a4paper,landscape,columns=3]{cheatsheet} %https://github.com/ACHinrichs/LaTeX-templates/blob/master/cheatsheet/cheatsheet.pdf

\usepackage[english]{babel}
\usepackage[utf8]{inputenc}
\usepackage[]{graphicx}
\usepackage{algpseudocode}
\usepackage{color}
\usepackage{hyperref}
\usepackage[algo2e,algoruled]{algorithm2e}

\title{Utopia cheat-sheet}

\definecolor{cppred}{rgb}{0.6,0,0} % for strings
\definecolor{cppping}{rgb}{1,0.7,0.7} % for numbers
\definecolor{cppgreen}{rgb}{0.25,0.5,0.35} % comments
\definecolor{cpppurple}{rgb}{0.5,0,0.35} % keywords
\definecolor{cppdocblue}{rgb}{0.25,0.35,0.75} % cppdoc

\lstset
{
language=C++,
basicstyle=\scriptsize\ttfamily,
numberstyle=\tiny\color{cpppink},
keywordstyle=\color{cpppurple}\bfseries,
stringstyle=\color{cppred},
commentstyle=\color{cppgreen},
}

\newcommand{\cppdouble}{{\color{cpppurple} double}}
\newcommand{\cppint}{{\color{cpppurple} int}}
\newcommand{\cppif}{{\color{cpppurple} if}}
\newcommand{\cppwhile}{{\color{cpppurple} while}}
\newcommand{\cppelse}{{\color{cpppurple} else}}


\begin{document}
% \maketitle
\href{https://bitbucket.org/zulianp/utopia}{\includegraphics[width=0.5\linewidth]{../../utopia_logo.pdf}}
\section{Installation}\footnotesize

\begin{lstlisting}[language=bash]
git clone --recurse-submodules  \
https://bitbucket.org/zulianp/utopia.git

cd utopia/utopia && mkdir build && cd build
cmake .. -DCMAKE_INSTALL_PREFIX=<your install folder>
make && make install
\end{lstlisting}

Click  \href{https://bitbucket.org/zulianp/utopia/wiki/Getting%20started}{HERE} for more details!

\section{Basics}
\subsection{Backends}\footnotesize
Each backend is defined by their tensor types. We denote vector types with V, Matrix types with M, general tensor types with T.
\begin{itemize}
    \item
    \textbf{Blas--Lapack}, node-level dense algebra
    \begin{lstlisting}
    using M = BlasMatrixd;
    using V = BlasVectord;
    \end{lstlisting}
    \item
    \textbf{PETSc}, parallel sparse and dense algebra
    \begin{lstlisting}
    using M = PetscMatrix;
    using V = PetscVector;
    \end{lstlisting}
    \item
    \textbf{Trilinos}, parallel sparse algebra
    \begin{lstlisting}
    using M = TpetraMatrixd;
    using V = TpetraVectord;
    \end{lstlisting}
\end{itemize}


\subsection{Traits}\footnotesize
Traits are used to access data-types associated with a tensors T.
\begin{lstlisting}
    // The MPI Communicator
    using Communicator = Traits<T>::Communicator;

    // Double or single precision real number
    using Scalar = Traits<T>::Scalar;

    // 32 or 64 bits integer number used for global indexing
    using SizeType = Traits<T>::SizeType;

    // 32 or 64 bits integer number used for local indexing
    using LocalSizeType = Traits<T>::LocalSizeType;

    // Array of indices
    using IndexArray = Traits<T>::IndexArray;

    // Array of scalars
    using ScalarArray = Traits<T>::ScalarArray;

    // Vector layout type specific to the backend
    using Layout = Traits<T>::Layout;

    // Matrix layout type specific to the backend
    using MatrixLayout = Traits<T>::MatrixLayout;
\end{lstlisting}

\section{Communicator}
\begin{lstlisting}
// Get the serial communicator
auto self = Communicator::self();

// Get the world communicator
auto world = Communicator::world();

// Get the default communicator (typically world)
auto comm = Communicator::get_default();
\end{lstlisting}

\section{Layouts}

\begin{lstlisting}
// Create a vector layout with n_local entries
// for each process and n entries in total
auto vl = layout(comm, n_local, n);

// Create a matrix layout from the vector layout
auto ml = square_matrix_layout(vl);

// Create a matrix layout with local and total
// number of rows and columns
auto ml = layout(comm, rows_local, cols_local, rows, cols);
\end{lstlisting}

\section{Tensor}
\subsection{Vector}
\begin{lstlisting}
// Construct a vector with a layout vl
V vec(vl);

// Construct a vector with and set it a uniform value
V vec(vl, val);

// Construct an existing vector with zeros,
vec.zeros(vl);
// or with a uniform value
vec.values(vl, val);

\end{lstlisting}
\subsection{Matrix}

\section{BLAS}
Dense and sparse basic linear algebra subroutines
\subsection{Level 1}

\subsection{Level 2}

\subsection{Level 3}

\subsection{Others}


\section{Linear solver}


\section{Nonlinear solver}

\section{Input-ouput}
\begin{lstlisting}
// Any tensor t can be inspected in the terminal
disp(t);
\end{lstlisting}


\section{Conversions and interoperability}


\section{Host-side manipulations}


\section{Device-side manipulations}


\section{Acknowledgements}
\cite{utopiagit}

\bibliographystyle{alpha}
\bibliography{Utopia_CheatSheet}

\end{document}